\scanpage{34}
(savoir le quotient par l'image de $\pi_2(Y)$),
et $=$ par $\pi_1(V^*_{z,Z})$ lui-$=$ si $\pi_2(Y)$ est
nul. Si on suppose pour simplifier
\deleted{\unsure{cyclique}}
que $\pi_2(Y)\longrightarrow\pi_1(V^*_{z,Z})$ est nul, on trouve
donc un diagramme qui précise (10),

\noindent (14)
\diagram{\begin{tikzcd}[column sep=small,row sep=small]
& \deleted{\unsure{\pi_1}} \\
1 \arrow[r,"{\deleted{\unsure{(14)}}}"'] & \pi_1(V^*_{z,Z}) \arrow[r] \arrow[d,equal] & \pi_1(V^*_{Y,X}) \arrow[r] \arrow[dd] \arrow[dr,equal,shorten <=20pt,shorten >=6pt] & \pi_1(Y) \arrow[r,dash] & 1 \\
& (I\ \text{``inertie''}) && (D\ \text{``décomposition''}) \\
&& \pi_1(X^*)
\end{tikzcd}},

où la ligne horizontale est exacte.

Un cas intéressant est celui où les
``pièces de construction'' qui \deleted{\unsure{sont}} interviennent,
savoir $Y$, $X^*$ et $V^*_{\unsure{Z,z|z,Z}}$, sont des espaces
``$K(\pi,1)$'', i.e.
\[
(15)\qquad \pi_i(Y),\ \pi_i(X^*),\ \pi_i(V^*_{\unsure{Z,z|z,Z}})\quad\text{nuls pour }i\geq 2,
\]
alors on en déduit qu'il en est de $=$
pour $V^*_{Y,X}$ (et bien sûr pour $V^*_{Y,X}\xleftarrow{\approx\,\mathrm{homot.}}Y$).
Si on a la ``théorie souple des \uline{liens} des
types d'homotopie'' à laquelle je faisais
\mspage{33}
